\thispagestyle{empty}

% TÍTULO DO RESUMO
\begin{center}
    \textbf{\large RESUMO}
\end{center}

\vspace{14pt}

% CORPO DO RESUMO EM PARÁGRAFO ÚNICO E ESPAÇAMENTO SIMPLES (ABNT NBR 6028)
\begin{singlespace}
\noindent\justifying
O presente trabalho apresenta o desenvolvimento do WeFIND, uma aplicação móvel colaborativa projetada para apoiar tutores e a comunidade na divulgação e na busca por animais perdidos e encontrados. A concepção do projeto baseou-se em uma pesquisa de campo e na análise comparativa de plataformas similares, que contribuíram para identificar dificuldades relacionadas à dispersão das informações, à localização dos casos e à exposição de dados de contato. A partir desse levantamento, foram definidos requisitos e modelos UML para orientar o desenvolvimento. A aplicação reúne recursos de geolocalização, consulta e filtragem de publicações, registro de avistamentos, comunicação privada e geração de cartazes com código QR para compartilhamento ou impressão. O desenvolvimento utiliza React Native e Expo no ambiente móvel e a plataforma Supabase com banco de dados relacional PostgreSQL para autenticação, armazenamento e sincronização de dados. Verificações funcionais manuais exploratórias e uma avaliação de usabilidade com usuários finais indicaram percepção positiva de facilidade de uso, consideradas as limitações da amostra e dos procedimentos adotados.

\vspace{18pt}

\noindent
\textbf{Palavras-chave:} Animais perdidos. Aplicativos móveis. Geolocalização. React Native. Supabase. Sistemas colaborativos.
\end{singlespace}

\clearpage
